%==============================================================
% BAB III  METODE
%==============================================================

\chapter{METODE}

Bab ini dapat diawali dengan kerangka pendekatan studi.
Metode penelitian dapat berupa percobaan laboratorium, percobaan lapangan,
dan survei lapangan yang dirancang sesuai dengan tujuan atau jenis penelitian,
seperti: eksploratif, deskriptif, koreksional, kausal, komparatif, eksperimen,
tindakan (\textit{action research}), pemodelan, analisis suatu teori, atau
kombinasi dari berbagai jenis penelitian tersebut.
Untuk penelitian yang menggunakan metode kualitatif, jelaskan pendekatan yang
digunakan, proses pengumpulan dan analisis informasi, dan proses penafsiran
hasil penelitian.
Maksud dari perincian ini ialah untuk menjamin keterulangan hasil.


%--------------------------------------------------------------
\section{Waktu dan Tempat}
%--------------------------------------------------------------

%--------------------------------------------------------------
\section{Alat dan Bahan}
%--------------------------------------------------------------

%--------------------------------------------------------------
\section{Prosedur Kerja}
%--------------------------------------------------------------

%--------------------------------------------------------------
\section{Analisis Data}
%--------------------------------------------------------------

